% Generated draft: requires mathematical and source review.
\documentclass[11pt]{article}
\usepackage[margin=25mm]{geometry}
\usepackage{fontspec}
\setmainfont{Latin Modern Roman}
\usepackage{amsmath,amssymb,mathtools,unicode-math}
\setmathfont{Latin Modern Math}
\usepackage{enumitem,needspace,xurl,hyperref}
\setlength{\parindent}{0pt}
\setlength{\parskip}{5pt}
\setlength{\emergencystretch}{3em}
\newcommand{\N}{\mathbb{N}}
\newcommand{\Z}{\mathbb{Z}}
\newcommand{\Q}{\mathbb{Q}}
\newcommand{\R}{\mathbb{R}}
\newcommand{\C}{\mathbb{C}}
\newcommand{\F}{\mathbb{F}}
\newcommand{\A}{\mathbb{A}}
\newcommand{\PP}{\mathbb{P}}
\newcommand{\E}{\mathbb{E}}
\newcommand{\eps}{\varepsilon}
\newcommand{\dd}{\,\mathrm{d}}
\begin{document}
\section*{UnsolvedMath ID: 30000532}
\textbf{Source problem code:} catalogue entry
\subsection{Definitions and mathematical statement}
\textbf{Two foundational questions.} What are motivic sheaves over \(k[t]/(t^2)\), and what is the tangent space to the space of motives? For the intended field \(k\), let \(R=k[\varepsilon]/(\varepsilon^2)\), the dual numbers. The map \(R\to k\), \(\varepsilon\mapsto0\), defines a first-order nilpotent thickening of \(\operatorname{Spec}k\). This nonreduced base, not merely its underlying reduced point, is part of the question.\par A possible precise organization first requires a chosen category or higher category \(\mathcal M(R)\) of the intended motivic sheaves, its counterpart \(\mathcal M(k)\), and a reduction functor \(i^*:\mathcal M(R)\to\mathcal M(k)\). For a fixed motive \(M_0\in\mathcal M(k)\), first-order deformations would be objects \(M_R\) equipped with an identification \(i^*M_R\cong M_0\), modulo the chosen notion of equivalence and any required flatness or other admissibility conditions. This describes the data a deformation problem must specify; it does not define the missing motivic category or assert that its deformation classes form a vector space.\par A tangent space must be taken at a specified point of a specified moduli problem. A set of isomorphism classes, a deformation groupoid, and the tangent complex of a derived moduli object are different possible outputs. No automatic identification with \(\operatorname{Ext}^1(M_0,M_0)\) is made without a category, enhancement, and deformation theorem. The source does not fix pure or mixed motives, integral or rational coefficients, a topology or realization for sheaves, the field characteristic, the moduli functor, or which infinitesimal information should be retained rather than discarded by passage to the reduced base. Thus these remain construction and definition questions, not a single stated isomorphism conjecture.\par The supplied triage checked 2026-08-21 labels the record open, describing later deformation-theoretic frameworks without a solution to the broad motivic formulation. Its framed-sheaf reference is not asserted here to define all motivic sheaves over dual numbers. This dated assessment requires independent verification.
\subsection{Short English statement}
Define the intended motivic objects over the dual numbers, then formulate their first-order deformations around a fixed motive. The category and moduli problem are missing, so a deformation set, groupoid, tangent space, and tangent complex cannot be identified without further definitions.
\subsection{Sources}
\begin{itemize}\item Local UnsolvedMath registry, permanent ID 30000532, code catalogue entry; \emph{Motivic Sheaves over Dual Numbers}, Oberwolfach Reports Open Problems.\item \emph{Algebraic K-Theory (2007)}, Oberwolfach Reports, DOI 10.4171/owr/2006/32, pp. 1941--1990. \url{https://doi.org/10.4171/owr/2006/32}. Proposed year and year tag: 2006; no proposer named.\item Supplied literature label \emph{OWR Algebraic K-Theory report}: \url{https://doi.org/10.4171/owr/2006/32}.\item Supplied literature label \emph{Framed-sheaf deformation theory}: \url{https://doi.org/10.1007/s00209-021-02802-x}.\item Literature triage checked 2026-08-21, status \texttt{open}; record updated 2026-09-16. The \texttt{corrected\_verified} field records notation and punctuation normalization. Verify dated claims independently.\end{itemize}
\end{document}
